\section{Tabelas de Frequência}

Um procedimento para entender o comportamento e distribuição dos dados consiste no uso
tabelas de frequência, especialmente para variáveis qualitativas, que não podem ser
resumidas por medidas de estatística descritiva a serem vistas mais adiante.
Considere um quadro de dados $D$ com $n$ registros e com um atributo $X$ com $k$ valores
possíveis iguais a $d_1, \ldots, d_k$.
Numa tabela de frequência para $D$ em relação ao atributo $X$ constam, geralmente,
duas colunas: a primeira é frequência (absoluta), i.e., o número de registros
$n_i$ tais que $X = d_i$, enquanto a terceira é a frequência relativa $f_i = n_i / n$,
também chamada de proporção e representada geralmente em porcentagem, denotando a fração
de regitros que pertecem à classe $d_i$.

As tabelas de frequência também podem ser usadas para descrever dados quantitativos
desde que esses sejam devidamente discretizados, o que ocorre, comumente, por meio
de intervalos de amplitude pré-determinada, de forma que cada intervalo passa a representar
uma classe.
Diferentes valores de amplitudades para cada classe podem levar a visualizações distintas.
Amplitudes pequenas levam a muitas classes, que por sua vez podem gerar perda de informação,
enquanto que amplitudes grandes criam poucas classes, prejudicando a tarefa de resumo.
O recomendado é a criação de 5 a 15 classes de mesma amplitude.

Para dados contínuos discretizados em intervalos de amplitude $\Delta_i$ para a classe
$d_i$, a tabela pode conter também: a densidade de frequência, dada por $n_i / \Delta_i$,
a densidade de proporção, obtida por $f_i / \Delta_i$, e a frequência acumulada, dada
por
\[
	\sum_{j = 1}^i f_i
\]
Essas novas colunas podem ser usadas para construir gráficos que permitam analisar
distribuição dos valores de uma variável numérica.

% Outra utilidade das tabelas de frequência é na construção da função de distribuição
% empírica.
% Suponha que $N_X(x)$ denote o número de observações para a variável $X$ tal que
% $X \le x$.
% Com isso, a função de distribuição empírica de $X$ é dada por
% \[
%     F_X(x) = \frac{N(x)}{n}
% \]
% onde $n$ é o numero de observações.
